% Generated draft: requires mathematical and source review.
\documentclass[11pt]{article}
\usepackage[margin=25mm]{geometry}
\usepackage{fontspec}
\setmainfont{Latin Modern Roman}
\usepackage{amsmath,amssymb,mathtools,unicode-math}
\setmathfont{Latin Modern Math}
\usepackage{enumitem,needspace,xurl,hyperref}
\setlength{\parindent}{0pt}
\setlength{\parskip}{5pt}
\setlength{\emergencystretch}{3em}
\newcommand{\N}{\mathbb{N}}
\newcommand{\Z}{\mathbb{Z}}
\newcommand{\Q}{\mathbb{Q}}
\newcommand{\R}{\mathbb{R}}
\newcommand{\C}{\mathbb{C}}
\newcommand{\F}{\mathbb{F}}
\newcommand{\A}{\mathbb{A}}
\newcommand{\PP}{\mathbb{P}}
\newcommand{\E}{\mathbb{E}}
\newcommand{\eps}{\varepsilon}
\newcommand{\dd}{\,\mathrm{d}}
\begin{document}
\section*{UnsolvedMath ID: 800228}
\textbf{Source problem code:} catalogue entry
\subsection{Definitions and mathematical statement}
Let \(D\subset\C^n\) be a domain, meaning a connected open set. Under the usual \(C^2\)-boundary convention, strong pseudoconvexity means that near each boundary point there is a real-valued defining function \(\rho\), with \(D=\{\rho<0\}\) locally and \(\dd\rho\ne0\) on the boundary, such that
\[
\sum_{j,k=1}^{n}\frac{\partial^2\rho}{\partial z_j\partial\overline z_k}(q)v_j\overline v_k>0
\]
for every boundary point \(q\) in that neighborhood and every nonzero complex tangent vector \(v\) satisfying \(\sum_{j=1}^{n}\frac{\partial\rho}{\partial z_j}(q)v_j=0\). This explains a standard convention; the supplied source does not specify boundary regularity or boundedness.

For a holomorphic function \(f\) on \(D\), define its underlying zero set by
\[
Z_D(f)=\{z\in D:f(z)=0\}.
\]
Write \(H^p(D)\) for the holomorphic Hardy space on \(D\), once the exponent \(p\) and the boundary norm convention have been fixed. For such fixed choices, define
\[
\mathcal Z_p(D)=\bigl\{Z_D(f): f\in H^p(D),\ f\not\equiv0\bigr\}.
\]
\textbf{Problem.} Give necessary and sufficient geometric conditions on a subset \(E\subset D\) for
\[
E\in\mathcal Z_p(D),
\]
when \(D\) is strongly pseudoconvex.

This is a formulation of the request using underlying zero sets of nonzero functions. The supplied statement does not specify the exponent or range of exponents, ambient dimension, boundedness, boundary regularity, Hardy-space convention, or whether zero multiplicities must be retained. If zero divisors are intended, the target must instead retain the vanishing multiplicities along their components. The exclusion of the identically zero function is also a convention requiring confirmation. Consequently, the displayed formulation leaves these source qualifications unresolved.
\subsection{Short English statement}
Characterize geometrically which sets in a strongly pseudoconvex domain occur as zero sets of holomorphic Hardy-space functions. The supplied statement leaves the exponent, domain hypotheses, Hardy-space convention, and treatment of multiplicities unspecified.
\subsection{Sources}
\begin{itemize}
\item Local UnsolvedMath registry, permanent ID \(800228\), source code \texttt{catalogue entry}, title: \emph{Several Complex Variables Problem List --- Problem 24}. The record attributes the problem to Buzzard, Fornaess, and Krantz and gives the proposed year as 1995.
\item Buzzard, Fornaess, and Krantz, \emph{The Several Complex Variables Problem List} (1995), Problem 24. Supplied source link: \url{https://arxiv.org/abs/math/9509202}. The record reports access on 2026-10-16 and extraction from arXiv LaTeX; these provenance claims have not been independently verified here.
\item The local record labels the problem \emph{open}. This is a supplied registry status, not an independently verified assessment of the literature.
\end{itemize}
\end{document}
